\documentclass[11pt]{article}

\usepackage[margin=1in]{geometry}
\usepackage{parskip}
\usepackage{amsmath, amssymb, amsthm}
\usepackage{booktabs}
\usepackage{hyperref}

\title{A small library for the binomial distribution}
\author{Group N}
\date{}

\begin{document}
\maketitle

\begin{abstract}
We present \texttt{binomial}, a small Python library for the binomial
distribution \(X \sim \text{Bin}(n, p)\). The library exposes five
short functions: the probability mass function, the cumulative
distribution function, the mean, the variance, and a simulator. The
binomial coefficient is computed exactly with
\texttt{scipy.special.comb}~\cite{virtanen2020}. We verify the
library on the fair-coin case \(n = 10\), \(p = 1/2\): the PMF and
CDF agree with the textbook values, the PMF sums to one, and the
empirical mean of \(5000\) simulated draws of \(\text{Bin}(20, 1/2)\)
lies within \(0.3\) of the theoretical mean of \(10\). The library
follows the Diataxis framework for documentation~\cite{procida2017}
and the combinatorics and probability chapters of the open-source
textbook \emph{Python for Mathematics}~\cite{knight2024}.
\end{abstract}

\section{Introduction}

The binomial distribution describes the number of successes in \(n\)
independent Bernoulli trials, each with success probability
\(p \in [0, 1]\). Writing \(X \sim \text{Bin}(n, p)\), the
probability mass function and the cumulative distribution function
are

\[
    P(X = k) = \binom{n}{k} p^k (1 - p)^{n - k},
    \qquad
    P(X \le k) = \sum_{j = 0}^{k} \binom{n}{j} p^j (1 - p)^{n - j},
\]

and the mean and variance are \(\mathbb{E}[X] = n p\) and
\(\text{Var}(X) = n p (1 - p)\)~\cite{ross2014}.

The binomial distribution is built into
\texttt{scipy.stats.binom}~\cite{virtanen2020}, which provides every
function we implement here. We wanted a small library that we could
read and understand line by line, with the binomial coefficient
computed exactly and the simulator written out as Bernoulli trials.

We present a five-function library,
\texttt{binomial}; we verify it on the fair-coin case \(n = 10\),
\(p = 1/2\); and we show that the empirical mean of \(5000\)
simulated draws of \(\text{Bin}(20, 1/2)\) is close to the
theoretical mean.

\section{The library}

The library lives in the single module \texttt{binomial.py} and
follows the combinatorics and probability chapters of \emph{Python
for Mathematics}~\cite{knight2024}.

\texttt{pmf(number\_of\_successes,
number\_of\_trials, success\_probability)} returns
\(\binom{n}{k} p^k (1 - p)^{n - k}\). The binomial coefficient is
computed exactly with \texttt{scipy.special.comb(\dots,
exact=True)}~\cite{virtanen2020}; we then multiply by the
probabilities.

\texttt{cdf(number\_of\_successes,
number\_of\_trials, success\_probability)} returns the partial sum
of the PMF up to \(k\). Summing eleven floats for a small \(n\) is
cheap, so we did not reach for \texttt{scipy.stats.binom.cdf}.

\texttt{mean(number\_of\_trials, success\_probability)} and
\texttt{variance(number\_of\_trials, success\_probability)} return
\(n p\) and \(n p (1 - p)\) respectively.

\texttt{simulate(number\_of\_trials, success\_probability,
number\_of\_simulations, seed)} draws each binomial sample as the
number of \(n\) independent uniform draws that fall below \(p\). The
``seed'' argument lets us reproduce a simulation.

\section{Worked example}

We illustrate the library on the fair-coin case \(n = 10\),
\(p = 1/2\). The PMF and CDF agree with the textbook values shown
in Table~\ref{tab:fair-coin}. The PMF sums to one and the CDF
reaches one at \(k = n\).

\begin{table}[h]
\centering
\begin{tabular}{r l l}
\toprule
\(k\) & \(P(X = k)\) & \(P(X \le k)\) \\
\midrule
0  & \(1/1024\)    & \(1/1024\) \\
1  & \(10/1024\)   & \(11/1024\) \\
2  & \(45/1024\)   & \(56/1024\) \\
3  & \(120/1024\)  & \(176/1024\) \\
4  & \(210/1024\)  & \(386/1024\) \\
5  & \(252/1024\)  & \(638/1024\) \\
\bottomrule
\end{tabular}
\caption{\textbf{The PMF and CDF of $X \sim \text{Bin}(10, 1/2)$.}
The mode is at \(k = 5\), with probability \(252/1024\). The library
returns these values to floating-point precision, with the binomial
coefficient computed exactly.}
\label{tab:fair-coin}
\end{table}

For \(\text{Bin}(20, 1/2)\) the theoretical mean is \(10\) and the
theoretical variance is \(5\). The empirical mean of \(5000\)
simulated draws lies within \(0.3\) of the theoretical mean in the
seeded run that we used as a test.

\section{Discussion}

\texttt{scipy.stats.binom} would have
covered everything. We wrote a small library because we wanted to
read the formulas inside, and because each function ends up being a
single line that maps to the definition in the probability
chapter~\cite{ross2014}.

The default
\texttt{scipy.special.comb} returns a float, which is fast but
already rounded for moderately large \(n\). With \texttt{exact=True}
the function returns a Python integer, so the factor in front of
\(p^k (1 - p)^{n - k}\) is exact and only the probabilities
contribute rounding error.

The library handles only the binomial
distribution. A natural extension is the multinomial distribution,
which generalises to more than two outcomes per trial. The same
approach (an exact multinomial coefficient and an explicit
simulator) would carry over.

\section{Conclusion}

We have presented \texttt{binomial}, a small library for the
binomial distribution. The library agrees with the textbook values
for the fair-coin case and the simulator produces samples whose
empirical mean is close to the theoretical mean.

\bibliographystyle{plain}
\bibliography{references}

\end{document}
